\subsection{Coevolution}
\label{sec:coevolution}

Evolutionary computation has drawn inspiration from co-evolutionary dynamics such as predation \citep{willkens2022evolving,popovici2012coevolutionary, ficici2000gametheoretic}, parasitism \citep{hillis1990coevolving,  cliff1995tracking}, competition \citep{stanley2004competitive}, and mutualism \citep{wang2019paired, bucci2005identifying, bucci2003mathematical, bucci2003focusing, watson2003computational, watson2001coevolutionary}.
However, the specific form of co-evolutionary relationship plays an often underappreciated role in these algorithms.

Under classic framing, evolutionary interactions are transient \citep{potter2000cooperative,ma2018survey, pan2022effective, gong2015distributed, goh2009competitivecooperative, antonio2017coevolutionary,hancer2025survey, wiegand2004analysis, garciapedrajas2005cooperative, wiegand2001empirical, byrski2015evolutionary}.
However, this classic framing is a clear case of non-symbiotic co-evolution, because each individual subcomponent is only guaranteed to be evaluated with another individual subcomponent once.
This structure is analogous to an individual hummingbird visiting an individual flower: they both benefit if the hummingbird's beak happens to perfectly fit the flower's shape, but the hummingbird's fitness is based on how well it matches the flowers on average (taking into account the hummingbird's ability to choose flowers of course) and the flower's fitness is based on how well it matches the visiting pollinators on average.
A life-long relationship between individuals of different species, i.e., a symbiotic relationship, guarantees that two individuals will continuously interact with each other, regardless of whether the relationship is parasitic or mutualistic.
These relationships can be highly asymmetric, with one organism being considered a \textit{host} and the other the \textit{symbiont}, in which case the ``life-long'' relationship is only
Mutualistic symbiosis frequently also involves vertical transmission, where the host offspring is infected with the symbiont's offspring near birth.
Vertical transmission is most often in the context of a host/symbiont system, i.e., when there is a strong asymmetry in the size and lifespan of the partners.
This mode of transmission enables the offspring of the partners to continue the same mutualistic relationship in terms of the shorter of the two lifespans.

For example, it has previously been shown that parasites can increase host behavioral complexity through the Population-Genetic Memory Hypothesis and Red Queen dynamics \citep{ladle1992parasites,zaman2014coevolution}.
% TODO It has also been hypothesized that the presence of both mutualistic and parasitic interactions in the same system could increase evolved host complexity beyond what the host would reach with only one type of interaction due to the more complex selection pressures from multiple interactions \citep{willkens2022evolving}.
Theory can also predict the parasitism-mutualism spectrum \citep{vostinar2019spatial,johnson2021coevolutionary}.

Parasite populations suppress fitness for genotypes in large sections of the fitness landscape, making the host populations adapt towards new alternate genotypes \citep{gupta2022hostparasite,zaman2014coevolution,kussell2006polymerpopulation, kumawat2021architecture}.
This host-parasite feedback has also been shown to aid evolvability in some cases, leading to the emergence of hosts with a genotype that changes phenotype with a single point mutation, and has a larger number of such point mutations that do change phenotype \citep{zaman2014coevolution}.

TODO cophylogenetic analyses?
